% ==============================================================================
% File: lpt_paper.tex
% Bài báo: The Surprising Effectiveness of Shared Memory in Looped Transformers
% Tác giả: Giovanni Monea, Keshav Ramji, Yousef El-Kurdi, Luis A. Lastras,
%          Yoav Artzi, Nathan Godey, Ramón Fernandez Astudillo
%          (Cornell và IBM Research, arXiv:2610.02383, 2026)
% Nguồn số liệu: reports/2610.02383_shared_memory_looped_transformers.md
% ==============================================================================

\section{Bộ nhớ dùng chung trong Transformer vòng lặp LPT}

% ------------------------------------------------------------------------------
% Slide 0: Tổng quan bài báo
% ------------------------------------------------------------------------------
\begin{frame}{LPT: Tổng quan bài báo}
    \begin{itemize}
        \item \textbf{Tên bài báo:} \textit{The Surprising Effectiveness of Shared Memory in Looped Transformers}
        \item \textbf{Tác giả và tổ chức:} Giovanni Monea, Keshav Ramji, Yousef El-Kurdi et al., Cornell và IBM Research, arXiv:2610.02383
        \item \textbf{Ý tưởng cốt lõi:} Huấn luyện Transformer vòng lặp từ đầu với bộ nhớ dùng chung: chỉ vòng lặp đầu ghi KV cache đầy đủ, các vòng sau đọc lại cache đó cộng một cửa sổ cục bộ ngắn.
        \item \textbf{Đóng góp nổi bật:}
        \begin{itemize}
            \item \textbf{Chia sẻ bộ nhớ không làm giảm chất lượng mà còn cải thiện nó:} LPT giảm PPL 0,38 đến 1,09 so với Transformer vòng lặp thường, bộ nhớ KV nhỏ hơn 2,0 đến 4,7 lần.
            \item \textbf{Mô hình lai H-LPT:} ở 5 vòng lặp, PPL thấp hơn Transformer chuẩn 1,12 đến 1,82 và dùng ít hơn 76 đến 79\% bộ nhớ ngữ cảnh.
            \item \textbf{Giải thích cơ chế:} đại lộ gradient, phân hóa vai trò giữa các vòng lặp.
        \end{itemize}
    \end{itemize}
\end{frame}

% ------------------------------------------------------------------------------
% Slide 1: Bối cảnh và thách thức
% ------------------------------------------------------------------------------
\begin{frame}{1. LPT: Bối cảnh và cái bẫy bộ nhớ KV cache}
    \textbf{Transformer vòng lặp:} áp dụng cùng một ngăn xếp $L$ tầng $R$ lần cho mỗi token.
    $$h^{(t)} = M_L\big(h_{\text{in}}^{(t)}\big), \qquad h_{\text{in}}^{(t+1)} = N_{\text{rec}}\big(h^{(t)}\big), \qquad t = 1, \dots, R$$
    \begin{itemize}
        \item Độ sâu hiệu dụng $L \cdot R$ nhưng số tham số chỉ bằng $L$ tầng.
    \end{itemize}
    \vspace{0.15cm}
    \textbf{Cái bẫy bộ nhớ:}
    \begin{itemize}
        \item Mỗi vòng lặp ghi một KV cache riêng, nên bộ nhớ KV tăng tuyến tính theo số vòng lặp: $M_{\text{KV}}^{\text{LT}} = L \cdot R \cdot T$.
        \item Ở ngữ cảnh dài, bộ nhớ ngữ cảnh có thể vượt bộ nhớ trọng số, phá hỏng lý do dùng mô hình vòng lặp.
    \end{itemize}
    \vspace{0.15cm}
    \textbf{Giải pháp cũ:} bỏ hoặc dùng lại cache lúc suy luận, không huấn luyện cho việc đó. Chỉ giữ cache của vòng đầu làm chất lượng của Ouro giảm mạnh. Chia sẻ bộ nhớ trước đây chỉ được xem là cách tiết kiệm, chấp nhận mất chất lượng.
\end{frame}

% ------------------------------------------------------------------------------
% Slide 2: Kiến trúc LPT
% ------------------------------------------------------------------------------
\begin{frame}{2. LPT: Kiến trúc Transformer vòng lặp dự đoán}
    \textbf{Tách việc xây bộ nhớ khỏi việc dự đoán theo trục vòng lặp:}
    \begin{itemize}
        \item \textbf{Vòng lặp trạng thái, $t=1$:} chạy như bình thường, ghi KV cache đầy đủ $C^{(1)}$ ở mọi tầng.
        \item \textbf{Vòng lặp dự đoán, $t \ge 2$:} đọc $C^{(1)}$ cộng cửa sổ cục bộ $\omega = 64$ vị trí gần nhất của chính nó, trong cùng một phép softmax:
    \end{itemize}
    $$\mathcal{A}_{\text{LPT}}^{(t)}(i) = C^{(1)}_{\le i} \;\cup\; C^{(t)}_{[i-\omega,\, i]}, \qquad t = 2, \dots, R$$
    \begin{itemize}
        \item \textbf{Bộ nhớ KV:} $M_{\text{KV}}^{\text{LPT}} = L\,[\,T + (R-1)\omega\,]$. Với $T = 4096$, $\omega = 64$, $R = 5$: nhỏ hơn mô hình vòng lặp thường 4,7 lần.
        \item \textbf{Mô hình lai H-LPT:} dùng nền mô hình lai Gated DeltaNet với một tầng chú ý đầy đủ trên năm tầng. Các tầng chú ý tuyến tính giữ trạng thái lặp riêng nhỏ cố định. Các tầng chú ý đầy đủ áp dụng quy tắc trên.
        \item \textbf{Huấn luyện:} mặt nạ khối của FlexAttention trên khóa giá trị dùng chung nối với khóa giá trị cục bộ.
    \end{itemize}
\end{frame}

% ------------------------------------------------------------------------------
% Slide 3: Thiết lập thực nghiệm
% ------------------------------------------------------------------------------
\begin{frame}{3. LPT: Thiết lập thực nghiệm}
    \begin{itemize}
        \item \textbf{Dữ liệu và cỡ mô hình:} FineWeb-Edu, ngữ cảnh 4096. Ba cỡ: S 152M, M 454M, L 1,01B tham số, huấn luyện 20B, 20B và 50B token.
        \item \textbf{Số vòng lặp:} 2, 3, 5. Phụ lục thêm 9.
        \item \textbf{Mô hình đối chứng:} Transformer chuẩn, Transformer sâu mà hẹp khớp độ sâu hiệu dụng và số tham số, Transformer vòng lặp thường, Transformer vòng lặp song song PLT, LOOPIE, mô hình lai LT2.
        \item \textbf{Hai hàm mục tiêu:} mất mát ở vòng cuối, và mất mát thích ứng của Ouro có phân phối dừng học được. Transformer vòng lặp thường dùng mất mát thích ứng, các mô hình còn lại dùng mất mát ở vòng cuối.
        \item \textbf{Đánh giá:} PPL trên 10M token giữ lại, độ chính xác zero-shot trung bình năm tác vụ (ARC-Easy, HellaSwag, PIQA, SciQ, LAMBADA), bộ nhớ và FLOPs tính giải tích tại ngữ cảnh 4096.
        \item \textbf{Không điều chỉnh siêu tham số:} dùng nguyên cấu hình mặc định của nanoGPT cho mọi mô hình.
    \end{itemize}
\end{frame}

% ------------------------------------------------------------------------------
% Slide 4: Kết quả chính
% ------------------------------------------------------------------------------
\begin{frame}{4. LPT: Kết quả chính ở cỡ 1B, 5 vòng lặp}
    \begin{table}[h]
        \centering
        \small
        \begin{tabular}{lccc}
            \toprule
            \textbf{Mô hình} & \textbf{PPL} $\downarrow$ & \textbf{Độ chính xác} $\uparrow$ & \textbf{Bộ nhớ KV so với chuẩn} $\downarrow$ \\
            \midrule
            Transformer chuẩn & 11,35 & 62,5 & $1{,}00\times$ \\
            Transformer vòng lặp thường & 11,39 & 63,4 & $5{,}00\times$ \\
            PLT & 10,59 & 66,4 & $1{,}06\times$ \\
            LT2 mô hình lai & 10,32 & 67,3 & $1{,}01\times$ \\
            \midrule
            \textbf{LPT} & \textbf{10,32} & \textbf{68,2} & $1{,}06\times$ \\
            \textbf{H-LPT} & \textbf{10,23} & 68,1 & $\mathbf{0{,}24\times}$ \\
            \bottomrule
        \end{tabular}
    \end{table}
    \vspace{-0.2cm}
    \begin{small}
    \begin{itemize}
        \item \textbf{LPT so với Transformer vòng lặp thường, mọi cỡ và độ sâu:} PPL thấp hơn 0,38 đến 1,09, độ chính xác cao hơn 0,4 đến 4,8 điểm, bộ nhớ KV nhỏ hơn 2,0 đến 4,7 lần, FLOPs tăng dưới 1\%.
        \item \textbf{H-LPT so với Transformer chuẩn, 5 vòng lặp:} PPL thấp hơn 1,12 đến 1,82, độ chính xác cao hơn 3,1 đến 5,6 điểm, bộ nhớ ngữ cảnh ít hơn 76 đến 79\%.
    \end{itemize}
    \end{small}
\end{frame}

% ------------------------------------------------------------------------------
% Slide 5: Độ sâu và độ ổn định
% ------------------------------------------------------------------------------
\begin{frame}{5. LPT: Mở rộng theo độ sâu và độ ổn định huấn luyện}
    \begin{itemize}
        \item \textbf{Transformer vòng lặp thường ngừng cải thiện từ 3 vòng lặp trở đi}, ở mọi cỡ. LPT tiếp tục cải thiện đến 5 vòng lặp.
        \item \textbf{Hàm mục tiêu ở vòng cuối:}
        \begin{itemize}
            \item Transformer vòng lặp thường huấn luyện bằng mất mát vòng cuối thì xấu dần khi tăng vòng lặp và sụp xuống mức unigram ở 5 vòng, với cấu hình chuẩn hóa của Ouro (S: PPL 2040).
            \item LPT ổn định, và mất mát vòng cuối tốt hơn mất mát thích ứng 0,16 đến 0,50 PPL: các vòng giữa được tự do tinh chỉnh trạng thái mà không phải dự đoán sớm.
        \end{itemize}
        \item \textbf{Chín vòng lặp:} từ 5 lên 9, LPT cải thiện nhẹ ở S nhưng giảm độ chính xác 1,7 điểm ở L. Mô hình lai H-LPT tiếp tục tốt lên ở mọi cỡ.
    \end{itemize}
    \begin{block}{Ý nghĩa}
        Chia sẻ bộ nhớ làm huấn luyện vòng lặp sâu ổn định hơn, nhưng không loại bỏ hoàn toàn bất ổn khi quá sâu.
    \end{block}
\end{frame}

% ------------------------------------------------------------------------------
% Slide 6: Cơ chế 1
% ------------------------------------------------------------------------------
\begin{frame}{6. LPT: Giải mã cơ chế 1 -- Đại lộ gradient}
    \small
    \textbf{Gradient từ mất mát cuối đến vòng lặp đầu yếu đi theo độ sâu:}
    \begin{itemize}
        \item Ở Transformer vòng lặp thường, chuẩn gradient đến vòng đầu giảm mạnh khi tăng số vòng lặp.
        \item Ở LPT, các vòng sau đọc trực tiếp KV cache của vòng đầu, nên gradient có đường tắt về vòng đầu. Ở 5 vòng lặp, phần lớn gradient đến vòng đầu đi qua đường này.
    \end{itemize}
    \vspace{0.1cm}
    \textbf{Bằng chứng nhân quả: cắt gradient của KV cache dùng chung.}
    \begin{itemize}
        \item Hướng thuận không đổi, chỉ cắt đường gradient. PPL tăng ở mọi cấu hình. Cỡ S, 5 vòng: 14,91 tăng lên 33,21. Cỡ M, 3 và 5 vòng: không huấn luyện được, PPL gần mức unigram.
    \end{itemize}
    \vspace{0.1cm}
    \textbf{Điểm cần thận trọng:}
    \begin{itemize}
        \item Tỉ lệ gradient qua KV dùng chung so với gradient qua trạng thái ẩn đã tăng theo độ sâu ngay từ lúc khởi tạo, nên xu hướng này tự nó chưa chứng minh lợi ích.
        \item Tác giả thừa nhận thí nghiệm chưa phân biệt được đường gradient ngắn hơn với thay đổi độ lớn hoặc hướng gradient.
    \end{itemize}
\end{frame}

% ------------------------------------------------------------------------------
% Slide 7: Cơ chế 2
% ------------------------------------------------------------------------------
\begin{frame}{7. LPT: Giải mã cơ chế 2 -- Phân hóa vai trò và cửa sổ cục bộ}
    \textbf{Các vòng dự đoán chú ý chủ yếu vào bộ nhớ dùng chung:}
    \begin{itemize}
        \item Tỉ lệ chú ý vào bộ nhớ dùng chung: 57 đến 71\% ở các vòng dự đoán đầu, 73 đến 80\% ở vòng cuối, kể cả khi cửa sổ cục bộ phủ toàn ngữ cảnh.
    \end{itemize}
    \vspace{0.1cm}
    \textbf{Khóa dùng chung và khóa cục bộ tách thành các cụm riêng trong PCA:}
    \begin{itemize}
        \item Ở Transformer vòng lặp thường, khóa của các vòng chồng lên nhau nhiều hơn.
        \item Vậy sự giống nhau giữa các vòng không phải điều kiện cần để chia sẻ bộ nhớ hiệu quả: vòng đầu chuyên xây bộ nhớ, các vòng sau chuyên dùng nó.
        \item Đây là bằng chứng định tính, trên một chuỗi kiểm tra.
    \end{itemize}
    \vspace{0.1cm}
    \textbf{Quét cửa sổ $\omega$ từ 0 đến 4096:}
    \begin{itemize}
        \item LPT tốt hơn Transformer vòng lặp thường ở mọi $\omega$. Cửa sổ ảnh hưởng ít, điều quyết định là được đọc bộ nhớ dùng chung. Chọn $\omega = 64$.
        \item Lưu ý: ở $\omega = 0$, vòng dự đoán vẫn chú ý vào mục nhập hiện tại của chính nó.
    \end{itemize}
\end{frame}

% ------------------------------------------------------------------------------
% Slide 8: Phép đo GPU
% ------------------------------------------------------------------------------
\begin{frame}{8. LPT: Đo thực tế trên GPU H100}
    \textbf{Bộ nhớ nhỏ hơn cho phép batch lớn hơn.} Ngữ cảnh 4096, bf16, một mã nguồn nghiên cứu dùng chung cho mọi mô hình.
    \begin{itemize}
        \item \textbf{Sức chứa batch:} LPT chứa 93 đến 98\% batch của Transformer chuẩn, và gấp 2,0; 2,9; 4,7 đến 4,8 lần Transformer vòng lặp thường ở 2, 3, 5 vòng lặp. H-LPT chứa gấp 3,6 đến 5,0 lần Transformer chuẩn.
        \item \textbf{Thông lượng của H-LPT so với Transformer chuẩn, tại batch tối đa:} 2,1 đến 2,2 lần ở 2 vòng; 1,4 đến 1,5 lần ở 3 vòng; 0,8 đến 0,9 lần ở 5 vòng.
    \end{itemize}
    \vspace{0.1cm}
    \textbf{Chi phí không biến mất:}
    \begin{itemize}
        \item Cùng batch, LPT chậm hơn Transformer vòng lặp thường 3 đến 8\%, vì vẫn phải đọc cả cache dùng chung. Bộ nhớ dùng chung tiết kiệm dung lượng, không tiết kiệm thời gian đọc.
        \item Batch bằng 1: một bước giải mã của Transformer vòng lặp thường chậm hơn chuẩn 2,3; 3,4; 5,7 lần ở 2, 3, 5 vòng. Ở 5 vòng LPT tốn 4,3 đến 4,8 lần FLOPs của Transformer chuẩn.
    \end{itemize}
\end{frame}

% ------------------------------------------------------------------------------
% Slide 9: Độ tin cậy và công bằng
% ------------------------------------------------------------------------------
\begin{frame}{9. LPT: Độ tin cậy thống kê và tính công bằng}
    \textbf{Bốn hạt giống ngẫu nhiên ở cỡ S, kiểm định t ghép cặp một phía:}
    \begin{table}[h]
        \centering
        \small
        \begin{tabular}{lccc}
            \toprule
            \textbf{Chênh lệch PPL} & \textbf{2 vòng} & \textbf{3 vòng} & \textbf{5 vòng} \\
            \midrule
            LPT so với Transformer vòng lặp thường & $-0{,}40$ & $-0{,}71$ & $-1{,}05$ \\
            H-LPT so với LT2 & $-0{,}04$ & $-0{,}07$ & $-0{,}13$ \\
            LPT so với PLT & $-0{,}07$ & $-0{,}27$ & $-0{,}52$ \\
            \bottomrule
        \end{tabular}
    \end{table}
    \vspace{-0.2cm}
    \begin{small}
    \begin{itemize}
        \item Độ lệch chuẩn giữa các hạt giống lớn nhất là 0,08 PPL. Lợi thế thật sự lớn chỉ khi từ 3 vòng trở lên. Ở 2 vòng, LPT so với PLT gần hòa.
        \item \textbf{Phép so sánh bị lệch có lợi cho LPT:} cấu hình chuẩn hóa của Ouro làm Transformer vòng lặp thường yếu. Thử thêm 5 cách chuẩn hóa khác, bản tốt nhất giảm PPL của nó 0,13 đến 0,76. Khi đó khoảng cách còn \textbf{0,31 đến 0,37} ở S và \textbf{0,06 đến 0,43} ở M.
        \item Chỉ cỡ S có nhiều hạt giống. Cỡ M và L là một lần chạy.
    \end{itemize}
    \end{small}
\end{frame}

% ------------------------------------------------------------------------------
% Slide 10: Hạn chế, phản biện và kết luận
% ------------------------------------------------------------------------------
\begin{frame}{10. LPT: Hạn chế, phản biện và kết luận}
    \textbf{Hạn chế:}
    \begin{itemize}
        \item Không tối ưu về tính toán: 5 vòng lặp tốn 4,3 đến 4,8 lần FLOPs giải mã. Đối chứng khớp tham số, không khớp FLOPs. Mọi token dùng cùng số vòng lặp.
        \item Tối đa 1B tham số, một kho dữ liệu, ngữ cảnh 4096, chưa thử ngữ cảnh dài. Chênh độ chính xác dưới 1 điểm nên xem là nhiễu.
        \item Chưa tối ưu siêu tham số, chưa công bố mã nguồn.
    \end{itemize}
    \vspace{0.1cm}
    \textbf{Kết luận:}
    \begin{itemize}
        \item Chia sẻ bộ nhớ là một thiên kiến quy nạp tốt cho mô hình vòng lặp, không chỉ là mẹo tiết kiệm bộ nhớ.
        \item H-LPT tạo biên chất lượng và bộ nhớ mới cho mô hình vòng lặp: PPL thấp hơn Transformer chuẩn với ít hơn 76 đến 79\% bộ nhớ ngữ cảnh.
        \item \textbf{Hướng tiếp theo:} số vòng lặp thích ứng theo token, nhiều vòng trạng thái hoặc xen kẽ vòng trạng thái và vòng dự đoán, ngữ cảnh dài.
    \end{itemize}
\end{frame}
